\section{代数的根}

设 A 是一个 k-代数。A 上的左伪模是一个 k-模 M，它带有伪环 A（II, 附录, No.~2, 第 378 页）上的左伪模结构，使得我们有 \(a(\lambda x) = \lambda(ax) = (\lambda a)x\)（对 \(\lambda \in k\)、\(a \in A\) 与 \(x \in M\)）。我们同样定义 A 上的右伪模。

设 M 是 A 上的一个左伪模；我们用下述方式在 M 上定义一个称为典范的左 \(\tilde{A}\)-模结构

\[
(\lambda e + a) x = \lambda x + a x \quad \text { for } \lambda \in k, a \in A, x \in M.
\]

反之，每个左 \(\widetilde{A}\)-模通过把算子环限制到 k 与 A，都典范地带有 A 上的一个左伪模结构。于是，A 上的左伪模的结构种与 \(\widetilde{A}\) 上的模的结构种是等价的（集合论, IV, §1, No.~7, 第 262 页）。

设 M 是 A 上的一个左伪模，并设 N 是 M 的一个 k-子模。则 N 是 M 的 \(\widetilde{A}\)-子模，当且仅当它在 A 的作用下稳定；这时我们称 N 是 M 的一个子伪模。

如同环的情形，我们定义 A 上的左伪模 \(A_{s}\) 与右伪模 \(A_{d}\)。A 的左（相应地，右）理想是 \(A_{s}\)（相应地，\(A_{d}\)）的子伪模。

设 \(\mathfrak{a}\) 是 A 的一个正则左理想，\(u\) 是模 \(\mathfrak{a}\) 的一个右单位元。置 \(\mathrm{M} = \mathrm{A}_s / \mathfrak{a}\)，并把 \(u\) 在 M 中的像记为 \(z\)。我们有 \(\mathrm{M} = \mathrm{A}z\)，且 \(\mathfrak{a}\) 是 \(z\) 的零化子。

反之，设 M 是 A 上的一个左伪模，并设 z 是 M 的一个元素，使得 M = Az。则存在 A 的一个元素 u，使得 z = uz。对每个 \(a \in A\)，我们有 \((au - a)z = 0\)，故 au − a 属于 z 的零化子 \(\mathfrak{a}\)。于是，\(\mathfrak{a}\) 是 A 的一个正则左理想，元素 u 是模 \(\mathfrak{a}\) 的右单位元，并且在过渡到商时，映射 \(a \mapsto az\) 定义一个从 \(A_{s}/\mathfrak{a}\) 到 M 的同构。

定义 2. —— 我们称 A 上的伪模 M 是单的，如果我们有 AM ≠ 0，并且 0 与 M 是 M 的仅有的一些子伪模。

这等价于说 \(\tilde{A}\)-模 M 是单的且其零化子不包含 A。当 A 是一个环而 M 是一个 A-模时，我们有 AM = M，故定义 2 与 VIII, 第 45 页的定义 1 一致。

命题 5. —— a) 设 \(\mathfrak{m}\) 是 A 的一个正则极大左理想。则伪模 \(A_s / \mathfrak{m}\) 是单的，并且存在 \(A_s / \mathfrak{m}\) 的一个非零元素，其零化子为 \(\mathfrak{m}\)。

\begin{enumerate}
\def\labelenumi{\alph{enumi})}
\setcounter{enumi}{1}
\item
  设 M 是一个单伪模，并设 x 是 M 的一个非零元素。我们有 M = Ax，x 的零化子 \(\mathfrak{m}\) 是 A 的一个正则极大左理想，并且在过渡到商时，映射 \(a \mapsto ax\) 定义一个从 \(A_{s}/\mathfrak{m}\) 到 M 的同构。
\item
  设 M 是一个不约化为 0 的伪模。则 M 是单的，当且仅当对 M 的每个非零元素 x 我们有 M = Ax。
\end{enumerate}

设 \(\mathfrak{m}\) 是 A 的一个正则极大左理想。我们已经看到，存在一个非零元素 \(z\)（\(\mathrm{A}_s / \mathfrak{m}\) 的元素），其零化子等于 \(\mathfrak{m}\) 且使得 \(\mathrm{A}z = \mathrm{A}_s / \mathfrak{m}\)；特别地，我们有 \(\mathrm{A}(\mathrm{A}_s / \mathfrak{m}) \neq 0\)。此外，\(\mathrm{A}_s / \mathfrak{m}\) 的每个子伪模都形如 \(\mathfrak{n} / \mathfrak{m}\)，其中 \(\mathfrak{n}\) 是 A 的一个包含 \(\mathfrak{m}\) 的左理想。由于 \(\mathfrak{m}\) 是极大的，仅有的可能是 \(\mathfrak{n} = \mathfrak{m}\) 与 \(\mathfrak{n} = \mathrm{A}\)，从而伪模 \(\mathrm{A}_s / \mathfrak{a}\) 是单的。这就证明了 a)。

在 b) 的假设下，M 中使得 Ay = 0 的元素 y 所成之集是 M 的一个异于 M 的子伪模，从而约化为 0。于是，Ax 是 M 的一个非零子伪模，由此推出 M = Ax。断言 b) 由定义 2 前的注记得出。

设 M 是一个非零伪模。假设对 M 的每个非零元素 x 我们有 Ax = M。特别地，我们有 AM ≠ 0。设 N 是 M 的一个非零子伪模，并设 x 是 N 的一个非零元素。我们有 Ax = M，因此 N = M。故 M 是单的。反之，若 M 是单的，则由 b)，对每个 \(x \neq 0\) 我们有 M = Ax。

定义 3. —— k-代数 A 的根，记作 \(\Re(\mathrm{A})\)，是 A 的正则极大左理想的交。

当 A 是一个环时，A 的每个左理想都是正则的，故根的定义与 VIII, 第 154 页的定义 2 一致。

例. —— 1) \(k\)-代数 \(\operatorname{End}_k^{\mathrm{f}}(\mathrm{V})\) 的根约化为 0（VIII, 第 436 页，例 1）。*代数 \(\mathcal{C}_0(\mathrm{T})\) 与 \(\mathrm{L}^1 (\mathbf{R})\)) 的根亦然。*

\begin{enumerate}
\def\labelenumi{\arabic{enumi})}
\setcounter{enumi}{1}
\tightlist
\item
  设 A 是一个所有元素都是幂零元素的伪环。A 的根等于 A，因为 A 是唯一的正则左理想。
\end{enumerate}

命题 5 立即给出下面的结果。

命题 6. —— 代数 A 的根是单伪模的零化子的交。它特别是 A 的一个双边理想。

命题 7. —— A 的根是 \(\widetilde{A}\) 的根在 A 上的迹；它也等于反代数 \(A^{\circ}\) 的根（即 A 的正则极大右理想的交）。若环 k 无根，则 A 的根等于 \(\widetilde{A}\) 的根。

等式 \(\Re (\widetilde{\mathrm{A}})\cap \mathrm{A} = \Re (\mathrm{A})\) 由 VIII, 第 438 页，命题 4, b) 得出。由于 \(\widetilde{\mathrm{A}}\) 与 \(\widetilde{\mathrm{A}}^{\circ}\) 有相同的根（VIII, 第 156 页，推论 1），等式 \(\Re (\mathrm{A}) = \Re (\mathrm{A}^{\circ})\) 得出。若 \(k\) 无根，则 \(\widetilde{\mathrm{A}}\) 的包含 A 的极大左理想的交等于 A。于是 \(\Re (\widetilde{\mathrm{A}})\) 包含于 A，从而等于 \(\Re (\mathrm{A})\)。

注. —— 设 x 与 y 是 A 的元素；我们称 x 是 y 的左拟逆元，或 y 是 x 的右拟逆元，如果在 \(\widetilde{A}\) 中，x − e 是 y − e 的左逆，即如果我们有 \(x + y = xy\)。由命题 7 与雅各布森定理（VIII, 第 156 页，定理 1），A 的根由 A 中这样的元素 x 组成：对 \(\widetilde{A}\) 中的每个 u，ux − e 在 \(\widetilde{A}\) 中左可逆。由于我们有 \((a + \lambda e)x - e = ax + \lambda x - e\)（对 \(a \in A\) 与 \(\lambda \in k\)），A 的根是 A 中这样的元素 x 所成之集：\(ax + \lambda x\)（对每个 \(a \in A\) 与每个 \(\lambda \in k\)）在 A 中有左拟逆元。

